
La validación técnica del sistema híbrido DEB+GP en el \cref{ch:cap4} demostró que el modelo alcanza precisión predictiva aceptable (MAPE = \num{8.3}~\%, R$^2$ = \num{0.78}) bajo condiciones controladas de laboratorio. Sin embargo, la precisión algorítmica no garantiza por sí sola la utilidad operativa: un sistema puede ser estadísticamente robusto y, al mismo tiempo, inútil para un productor agroindustrial que carece de formación técnica en modelado matemático o ciencia de datos \parencite{kamilarisReviewPracticeBig2017}. Esta sección define los criterios y métricas mediante los cuales se evalúa que la plataforma de soporte a la decisión —y no solo sus componentes predictivos— cumple efectivamente su función de traducción entre predicción y acción en condiciones de campo.

La validación de utilidad se distingue de la validación de precisión en tres dimensiones fundamentales: (i)~la \textit{interpretabilidad} de los indicadores visuales para operadores no expertos; (ii)~la \textit{latencia} de la información respecto a la dinámica del proceso biológico; y (iii)~la \textit{pertinencia contextual} de las recomendaciones operativas en el entorno real de pequeños productores del sur global.

\subsection{Métricas de usabilidad para operadores no expertos}
\label{subsec:cap5_usabilidad}

La interfaz de la plataforma se evalúa mediante un protocolo de prueba de usuario estructurado, aplicado a un grupo de operarios de plantas de bioconversión con nivel educativo diverso (bachillerato a técnico agropecuario). Las métricas de usabilidad se definen como:

\begin{itemize}
    \item \textbf{Tasa de interpretación correcta del semáforo:} porcentaje de participantes que, ante un escenario simulado de nivel rojo, identifican correctamente la acción requerida (volteo inmediato) en menos de \num{10}~segundos de exposición al dashboard. El criterio de aceptación es $>$ \num{90}~\%, dado que el semáforo condensa la complejidad del modelo en un único indicador cromático de alto contraste.
    
    \item \textbf{Tiempo de localización de información ($\tau_{\text{loc}}$):} tiempo transcurrido entre la solicitud de encontrar un dato específico (por ejemplo, ``¿Cuánto carbono se emitió ayer?'') y la identificación correcta del valor en la interfaz. El objetivo es $\tau_{\text{loc}} <$ \num{30}~s, garantizando que el operador no requiera navegación profunda ni conocimiento de estructuras de base de datos.
    
    \item \textbf{Tasa de error de acción ($\epsilon_{\text{acc}}$):} frecuencia con la que un participante ejecuta una acción distinta a la recomendada por el sistema ante un escenario dado (por ejemplo, aplicar riego en lugar de volteo ante nivel rojo). Se busca $\epsilon_{\text{acc}} <$ \num{10}~\%, validando que las instrucciones textuales del protocolo son inequívocas.
\end{itemize}

Estas métricas se complementan con una escala de carga cognitiva subjetiva (NASA-TLX adaptada) aplicada tras la interacción con el dashboard, evaluando la demanda mental, el esfuerzo físico percibido y la frustración del operador. Un sistema de soporte a la decisión efectivo debe reducir, no aumentar, la carga cognitiva respecto a la gestión empírica tradicional basada en inspección visual olfativa del sustrato.

\subsection{Latencia end-to-end y pruebas de conectividad rural}
\label{subsec:cap5_latencia}

El requisito de latencia $<$ \num{1}~minuto entre la lectura del sensor NDIR en el nodo edge y la visualización en el dashboard del operador se valida mediante pruebas de campo en tres escenarios de conectividad representativos del sur global:

\begin{enumerate}
    \item \textbf{Escenario A — Conectividad estable (Wi-Fi local):} el nodo transmite vía MQTT sobre red Wi-Fi interna del galpón. La latencia medida incluye: adquisición del sensor (\num{1}~s), filtrado edge (\num{0.5}~s), transmisión MQTT (\num{0.3}~s), ingesta backend (\num{0.2}~s), ejecución motor DEB+GP (\num{5}~s para inferencia GP + integración EDO), y renderizado dashboard (\num{0.5}~s). Total esperado: $\approx$ \num{7.5}~s, ampliamente dentro del umbral de \num{1}~min.
    
    \item \textbf{Escenario B — Conectividad intermitente (LoRaWAN):} el nodo almacena datos localmente durante desconexiones de hasta \num{72}~h y transmite en ráfaga cuando recupera señal. En este modo, la latencia de la información histórica es irrelevante para el control en tiempo real, pero el procesamiento edge local debe mantener el semáforo funcional. Se valida que el algoritmo de detección de \ce{CH4} $>$ \num{100}~ppm se ejecute en el ESP32 con latencia $<$ \num{5}~s desde la lectura, activando una alarma local (buzzer o LED) independientemente de la conectividad WAN.
    
    \item \textbf{Escenario C — Sin conectividad (modo autónomo):} el nodo opera completamente aislado durante \num{48}~h. La validación consiste en verificar que: (i)~el buffer circular no pierde datos; (ii)~el semáforo local (LED RGB en el gabinete del nodo) refleja correctamente el estado del reactor; y (iii)~al restaurar la conexión, la sincronización de datos rezagados no genera inconsistencias en el inventario dinámico de carbono (la integración numérica debe ser conmutativa respecto al orden de llegada de los paquetes).
\end{enumerate}

Los resultados de estas pruebas se reportan mediante una matriz de cumplimiento que contrasta el requisito contra el valor medido en cada escenario, incluyendo desviaciones atribuibles a congestión de red, interferencia electromagnética en ambientes húmedos o saturación del buffer por eventos de alta frecuencia (por ejemplo, volteo que genera picos transitorios de \ce{CO2} cercanos a \num{6000}~ppm).

\subsection{Pertinencia contextual y adopción por pequeños productores}
\label{subsec:cap5_pertinencia}

La utilidad final del sistema se mide por su capacidad de integrarse en la lógica productiva real de pequeños y medianos productores agroindustriales. Esta dimensión se evalúa mediante indicadores cualitativos y cuantitativos de adopción:

\begin{itemize}
    \item \textbf{Costo marginal del sistema:} el desembolso total por reactor (nodo sensor + licencia de software + mantenimiento anual) debe representar $<$ \num{5}~\% del valor del producto generado (biomasa seca + frass) en un ciclo estándar de \num{15}~días. Este umbral se deriva de estudios de adopción tecnológica en agricultura de pequeña escala, donde inversiones superiores al \num{10}~\% del ingreso operativo generan resistencia al cambio independientemente del beneficio técnico \parencite{fawzyStrategiesMitigationClimate2020}.
    
    \item \textbf{Mantenimiento local:} la capacidad del productor para reemplazar el sensor NDIR (\num{6000}~ppm) o el módulo de comunicación sin asistencia técnica externa. Se valida mediante un protocolo de ``desconexión-reconexión'' en campo, donde el operador intercambia el componente siguiendo un manual ilustrado de máximo \num{4} pasos. El criterio de éxito es la recuperación funcional del nodo en $<$ \num{20}~min sin intervención remota.
    
    \item \textbf{Soberanía de datos percibida:} encuesta de percepción aplicada al productor sobre quién posee, almacena y puede acceder a los datos de emisiones de su planta. El diseño de la plataforma prioriza el almacenamiento local (PostgreSQL en servidor de granja o mini-PC) sobre la nube centralizada, y la encuesta valida que esta arquitectura se traduce en una percepción de control efectivo por parte del usuario.
\end{itemize}

En conjunto, estas métricas de utilidad operativa constituyen la validación de que el sistema híbrido no es solo un artefacto de laboratorio, sino una herramienta de gestión ambiental integrable en las prácticas productivas reales del sur global. El cumplimiento de los umbrales definidos en esta sección habilita la transición del prototipo experimental (TRL~5) hacia una tecnología validada en entorno operativo (TRL~6), puente necesario para la futura escalabilidad comercial y el acceso a mecanismos de financiación climática.
